\documentclass[11pt]{article}
\usepackage{manual_docs_style}
\externaldocument{build/dataset_manifest_schema}[dataset_manifest_schema.pdf]
\externaldocument{build/candidate_generation}[candidate_generation.pdf]
\externaldocument{build/cheap_filter_metrics_stage}[cheap_filter_metrics_stage.pdf]
\externaldocument{build/surrogate_confidence_scoring_stage}[surrogate_confidence_scoring_stage.pdf]

\begin{document}

\docTitle{Stage: Generate Question}
\phantomsection\label{docstart:generate_question_stage}

This stage exports selected simulator runs into the canonical TSENV question bundle: the standard exchange format consumed by agentic rollout, baseline evaluation, and downstream analysis.

\section*{Required Inputs}

For each \termSimulator, question generation requires:

\begin{itemize}
\item \code{datasets/<dataset_id>/sample_manifest.json}: selected train/test panels.
\item \code{datasets/<dataset_id>/selected_runs.jsonl}: selected runs and split metadata.
\item \code{datasets/<dataset_id>/selection_report.json}: frozen selection diagnostics.
\item \code{runs/<run_id>/data.parquet}: materialized time-series artifacts.
\item \code{runs/model_record.json}: runtime inventory for selected runs.
\item \code{plans/<production_policy_id>/run_nodes.jsonl}: canonical production run recipes and hashes.
\item \code{plans/<production_policy_id>/run_edges.jsonl}: production baseline/intervention/time-zero relationships.
\item \code{experiment_config.json}
\item \code{description_levels.json}
\item \code{noise_adder.py}
\end{itemize}

Question generation itself should not need to rescore the whole candidate pool.
The preferred design writes the selected dataset artifacts first.

\section*{Step 1: Dataset Selection}

The selection step builds a selected dataset manifest.
It is documented in the Dataset Manifest Schema and is conceptually upstream of this stage.

Recommended invocation:

\begin{DocCode}
env/bin/python workflows/generate/select_dataset.py \
  --model SIMULATOR \
  --policy PRODUCTION_POLICY_ID \
  --dataset DATASET_ID \
  --surrogate-id SURROGATE_ID \
  [--row-slug ROW_SLUG [ROW_SLUG ...]]
\end{DocCode}

\section*{Step 2: Generate Exported Questions}

This step is deterministic once \code{sample_manifest.json} exists.

\begin{DocCode}
env/bin/python workflows/generate/create_exam_questions_tsenv_cls_from_registry.py \
  --model SIMULATOR [SIMULATOR ...] \
  --dataset DATASET_ID \
  --row-slug SLUG [SLUG ...]
\end{DocCode}

It materializes the canonical exported bundle as:
\begin{DocCode}
  tsENV_questions/<simulator_id>/
    questions.json
    dataframes/
      `-- *.parquet
    noise_adder.py
    sample_manifest.json
    model_record.json
\end{DocCode}
Where:
\section*{\texorpdfstring{\titlecode{sample_manifest.json}}{sample\_manifest.json}}

The manifest is keyed by \code{shot_slug}.
Each \code{shot_slug} maps to a non-empty list of entries, one per resolved \code{(shot_slug, test_set_slug, seed)} combination.

Each entry contains:

\begin{itemize}
\item \code{train_samples_baselines}
\item \code{train_samples}
\item \code{train_samples_hashes}
\item \code{test_samples_baselines}
\item \code{test_samples}
\item \code{test_samples_hashes}
\item \code{other_samples}
\item \code{seed}
\item \code{test_set_slug}
\item \code{shot_slug_recipe}
\item \code{accuracy_with_baselines}, when adversarial selection is evaluated
\item \code{accuracy_with_baselines_all_samples}, when evaluated
\item \code{train_test_sample_hash}
\end{itemize}

The pair \code{(test_set_slug, seed)} has a one-to-one mapping to \code{test_samples_hashes}.

\section*{\texorpdfstring{\titlecode{questions.json}}{questions.json}}

The exported \code{questions.json} contains:

\begin{itemize}
\item \code{version}
\item \code{questions}
\item \code{label_int_mapping}
\item \code{ground_truth_information}
\item \code{environment_name}
\end{itemize}

\code{ground_truth_information} is derived from the selected production run graph and selected dataset manifest.
For each selected sample, it records initial parameters, changed parameter, intervention time, new value, and \code{first_diff} values.
No-intervention samples use the no-intervention class and null intervention fields.
Surrogate confidence values may be copied into hidden provenance metadata, but they must not appear in the agent-visible prompt or files.

\section*{Question Text}

\code{question_text} contains the rendered prompt components:

\begin{itemize}
\item sample-source description;
\item environment description controlled by \termDescLevel;
\item observed-column descriptions;
\item intervention semantics;
\item allowed labels and label space;
\item no-change guidance;
\item few-shot context;
\item task artifact instructions;
\item prediction format;
\item mode-specific requirements;
\item evaluation and runtime constraints.
\end{itemize}

These values are resolved from \code{shared/config/tsenv_shared_description.json}, \code{description_levels.json}, and the selected dataset manifest.

\section*{Validation Rules}

\begin{itemize}
\item Question generation must not select new samples.
\item It may only export samples already present in \code{datasets/<dataset_id>/selected_runs.jsonl} and grouped by \code{sample_manifest.json}.
\item Every selected \code{run_id} must exist in \code{run_nodes.jsonl} and \code{model_record.json}.
\item Every selected run must have a materialized \code{data.parquet} artifact.
\item Test-set length must match the requested \code{number_test_samples}.
\item Test-set class counts should be balanced, with any two class counts differing by at most one.
\item Train-set length must equal \code{number_train_samples_per_class * num_labels}.
\item Train and test samples must not overlap within a question.
\item Comparable rows sharing \code{test_set_slug} and seed must share the same underlying test set.
\item Exported question hashes must be derived from selected sample hashes.
\end{itemize}

\end{document}
